\subsection{A smaller power base in the factorial-free alternative}
A separate companion starts from the incoming Factorial-free binomial CRT
source rather than from the shared-factorial-error source. Its retained
prime-power separation theorem uses
\[
 R\ge\operatorname{lcm}(1,\ldots,t)^2,\qquad
 B=\binom{R+t}{t},
\]
to obtain a suitable full prime-power factor for every row, and to assign
every prime-power factor of $B$ to a row for reverse CRT completeness.
Its product-label recovery originally required a larger evaluation point.
The following elementary refinement permits a simpler authentic power
request without changing its local row rules.

\begin{lemma}\label{lem:power-label}
Let $t\ge2$ be an integer and $F=t^t$. Suppose two multisets of labels
from $\{1,\ldots,t\}$ have equal products of weights $1+jF$, and at
least one multiset has cardinality at most $t$. Then the multisets agree,
including multiplicities.
\end{lemma}
\begin{proof}
First the cardinalities must agree. If the smaller count is $d\le t$
and the larger count is at least $d+1$, the larger product is at least
$(F+1)^{d+1}$, whereas the smaller is at most
\[
 (1+tF)^d\le t^d(F+1)^d.
\]
This contradicts $F+1>t^t\ge t^d$, including $d=0$.
Thus both cardinalities are the same $d\le t$.

Let $A(z),B(z)$ be the corresponding products of $d$ factors $1+jz$.
When $d=0$, both polynomials are one. Otherwise every coefficient from
degree zero through $d$ is positive, and for $0\le k\le d$ it is at most
\[
 \binom dk t^k=\binom d{d-k}t^k
       \le d^{d-k}t^k\le t^d\le F.
\]
Consequently every coefficient of $A-B$ has absolute value strictly less
than $F$: both coefficients being subtracted lie in $\{1,\ldots,F\}$.
A strict upper bound on either coefficient separately is unnecessary.
If $A-B$ were nonzero, divide it by its lowest nonzero power of $z$.
The resulting nonzero constant coefficient has magnitude less than $F$,
so evaluation at $F$ cannot vanish modulo $F$. This contradicts
$A(F)=B(F)$. Hence $A=B$. The distinct labels give distinct roots
$-1/j$, and unique factorization over $\Q$ recovers their multiplicities.
Labels zero, if allowed by the arithmetic convention, can first be
discarded because their weight is one.
\end{proof}

For this companion choose the cofinal even horizon
\begin{equation}\label{eq:power-horizon}
 t=16+2h_0^2,\qquad F=t^t,
\end{equation}
using the already present genuine power request at index $t$ with base
$t$. These are the meanings of existing quadratic power certificates,
not variable-exponent gates in the final polynomial. The base is an
existing register, so its expression no longer needs the incoming
$2(t+1)$ construction. The literal delta consists of two raw rows: $t_0=h_0^2$ becomes
$t_0=2h_0^2$, and the growth base $2t_0+34$ becomes $t_0+16$.
Every horizon use inherits its changed value through the same shared
$t_0$ register. The rebuilt frame includes the actual new base and every
other required request atom; the independent source-delta checker verifies
that complete propagation.

We check all new arithmetic compatibility conditions. For $t\ge16$,
$F=t^t>t!$, $F\ge16^{16}>2^{32}$ and
$v_2(F)=t\,v_2(t)\ge t$, because $t$ is even.
Taking $R=F^2$ therefore exceeds the retained CRT threshold.
For $1\le i\le t$, the inequality
$v_2(F^2)>\log_2t\ge v_2(i)$ gives
$v_2(F^2+i)=v_2(i)$. Hence
\[
 v_2\!\left(\binom{F^2+t}{t}\right)
 =\sum_{i=1}^t v_2(F^2+i)-v_2(t!)=0,
\]
so $B$ is odd, as required by the source's representative parity choices.

In completeness the retained beta spacing has the form
$y=F(2w+1)^2$. Choose the odd factor divisible by the odd part of $t!$
and arbitrarily large within that odd divisibility class. Since
$v_2(t!)<t\le v_2(F)$, this makes $y$ an arbitrarily large even
multiple of $t!$. The original beta CRT, even-code and odd-square radix
construction therefore remains available. The new horizon is cofinal,
and the retained padding theorem extends any finite proof to an admissible
horizon. The original exclusion of the boundary $y=F$, using the zero
initial accumulator and the terminal accumulator one, is unchanged.
All other quantitative arguments use only the now verified bounds
$F>t!$, $F\ge7$, $F>2^{32}$ and the retained endpoint inequalities
$x\ge F+1$ and $x\ge2+(t+1)y$.

The power companion also widens its eleven-remainder norm from
$M^2=1+4\sum R_i^2+G$ to $M^2=2\sum R_i^2+G$.
Here the retained odd-square radix gives $q\ge1$ independently.
The authenticated positive $F=t^t$ gives $y\ge F$, and the numerator
requests give nonnegative actual representatives and $H\ge2$.
The raw encoded-index row excludes $B=0$. The self-radix request then
forces $x\ge0$; the last beta endpoint excludes $x=0$, and the first
excludes a negative initial quotient. Hence
$x\ge F+1$ and $x\ge2+(t+1)y$, and
$q=|BH|^x\ge2^x>1$. Thus $M=q^x\ne0$ before using the new sphere.
The same coefficient estimates as in the retained source, using
$F\ge t$, $F\ge7$ and these endpoint bounds, place every canonical
low remainder below $L_i/32$. Dividing each normalized row by its
nonzero ratio gives an integral $r_i$ with
$|r_i|\le L_i/\sqrt2<3L_i/4$. The divisible difference from the true
remainder has magnitude less than $25L_i/32<L_i$, so it vanishes.
The original bound $|B|<q/2$ and the exact factorial-binomial
coefficient bound $C_*<q/32$ still give $B=C_*>0$ without circularity.
For completeness, the canonical remainders give
$2\sum R_i^2<22M^2/1024<M^2$, so their nonnegative integer gap has a
$G$ representation. This proves both directions of the additional sphere
change for this separate source.

For soundness, the unchanged signed bootstrap authenticates the request
$F=t^t$, so the retained prime-power theorem applies at $R=F^2$ and
Lemma~\ref{lem:power-label} supplies the product-label injectivity.
The unchanged local calculus and calibration recover a finite refutation.
Conversely, the cofinal and parity construction just checked supplies all
new power-base and horizon witnesses, and the retained reverse CRT theorem
supplies all other witnesses. Thus this companion has the same existence
interpretation, but by its own complete source route. It is not claimed
to be a coordinate bijection with the shared-error source, and it does
not impose $F=t!$. No least-possible evaluation point or global
coefficient optimum is asserted by Lemma~\ref{lem:power-label}.
